% ==================== 4. Experiments ====================
\section{Experiments}
\label{sec:exp}

\subsection{Dataset and protocol}
\label{sec:data}

\textbf{Capture.} A tripod fixes the camera; a phone shoots through a tinted car window. For
each scene a reflection-contaminated sequence is recorded, then the camera side is covered
with a black velvet cloth and the transmission reference is recorded. We shoot 13 training
and 12 frozen test scenes (631 pairs in total), at four controlled attenuation levels
(original plus $8/11/14\%$ manual darkening), and annotate, per frame, the occlusion region
and a signed reflection layer. Table~\ref{tab:data} summarizes the statistics and the two
acquisition defects that we disclose rather than hide.

\begin{table}[t]
\centering
\caption{Tinted-glass benchmark. Defects are reported, not corrected.}
\label{tab:data}
\small
\begin{tabular}{@{}p{1.9cm}p{5.6cm}@{}}
\toprule
Item & Value \\
\midrule
Pairs / scenes & 631 pairs; 13 training scenes \\
Resolution & test panel $1280\times832$; pools $955\times934$ (apr/may/jun), $1392\times774$ (re-crop) \\
Annotations & $I$, $T$, signed $R$ (uint8$+128$), per-frame $\Omega$ \\
Attenuation levels & 4 ($\rho$: apr $0.969$, may $0.899$, jun $0.921$, test $0.858$) \\
Occlusion present? & apr \ding{51}, jun \ding{51}, may \ding{55} (confirmed by inspection) \\
$\Omega$ area (train pool) & $15.3\%$ (82 new pairs), $20.8\%$ (126 old pairs) \\
Defect 1 & one GT per scene, median $374$ frames from the input (6\% within 50) \\
Defect 2 & inside $\Omega$ the input is $13\%$ darker than its own target ($I/T|_{\mathrm{in}}\approx0.87$): a per-region gain is required \\
\bottomrule
\end{tabular}
\end{table}

\textbf{Metrics.} PSNR, SSIM, HF-PSNR and a gradient-cosine term on 15 native-resolution
scenes; the minimum detectable effect of the paired panel is $0.26$~dB, and for
brightness-only effects $1.19$--$1.31$~dB.

\subsection{Zero-shot benchmark}
\label{sec:bench}

Table~\ref{tab:bench} evaluates seven public methods without any training on our data. Two
observations matter for the rest of the paper. First, the ranking tracks \emph{whether the
training data is real} rather than publication year: the real-pair method of
2024~\cite{zhu2024rrw} is first ($20.50$~dB) while synthetic-training methods fall to
$18.12$--$18.90$~dB, i.e.\ \textbf{below the identity baseline} ($18.44$~dB). Second, the
hole/non-hole gap is negative for every method, but \emph{its sign depends on the partition
caliber}: switching from an Otsu-based reflection mask to the training-time $\Omega$
criterion flips it for all methods. We therefore never use the raw gap as evidence; the
comparable quantity is the gain relative to identity inside the same region, which is at
most $+0.56$~dB inside $\Omega$ for every method (versus $+0.5$ to $+2.9$~dB outside).

\begin{table*}[t]
\centering
\caption{Gain over the trivial baseline (identity) in the two partitions of the frozen panel.
The partition used for evaluation is the v13 criterion --- a \emph{reflection} mask
($r=R/T$ gate $+$ Otsu, median $9.4\%$ of the frame) --- because it is the only criterion that
stays well behaved on these outdoor scenes; the training-time occlusion criterion is not used
for evaluation (see the text). Under this partition every method gains \emph{more} inside than
outside, so the raw negative hole/non-hole gap is a property of the region, not of the method.}
\label{tab:regiongain}
\small
\setlength{\tabcolsep}{6pt}
\begin{tabular}{@{}lccccc@{}}
\toprule
Method & PSNR & hole & gain in hole & non-hole & gain in non-hole \\
\midrule
identity (no processing) & 18.44 & 15.17 & $0$ & 18.97 & $0$ \\
official RDNet~\cite{zhao2025rdnet} & 20.05 & 19.25 & $+4.08$ & 20.19 & $+1.22$ \\
RRW~\cite{zhu2024rrw} & 20.50 & 19.21 & $+4.05$ & 20.68 & $+1.71$ \\
DSRNet~\cite{hu2023dsrnet} & 20.25 & 18.30 & $+3.14$ & 20.53 & $+1.56$ \\
LightRR & 19.18 & 16.86 & $+1.70$ & 19.51 & $+0.54$ \\
ERRNet~\cite{wen2019errnet} & 18.90 & 16.10 & $+0.94$ & 19.33 & $+0.36$ \\
IBCLN~\cite{li2020ibcln} & 18.12 & 15.48 & $+0.31$ & 18.52 & $-0.45$ \\
\bottomrule
\end{tabular}
\vspace{2pt}
\footnotesize Identity itself scores $-3.81$~dB in this partition, i.e.\ the region is
intrinsically harder to score; comparing \emph{gains} removes that. This partition is
\textbf{not} the occlusion region: it marks where reflection is \emph{weak}, while the occlusion
region marks where reflection was \emph{blocked} --- the two are close to opposite, and the
training-time criterion degenerates to a $35\%$ field on these scenes. No per-region number in
this paper is therefore reported under the occlusion criterion on the test panel.
\end{table*}

\textbf{What the partition can and cannot say.} Table~\ref{tab:regiongain} reports gains
relative to identity in both partitions. Two facts follow. First, \emph{every} method gains more
inside the v13 hole ($+0.3$ to $+4.1$~dB) than outside it ($+0.5$ to $+1.7$~dB), so the plain
negative hole/non-hole gap of Table~\ref{tab:bench} is a property of the region --- identity has
the same gap ($-3.81$~dB) --- and is not evidence that methods fail there. Second, this
partition is a \emph{reflection} mask, not the occlusion region, and the occlusion criterion we
use at training time degenerates on these outdoor scenes (median area $35\%$, against
$15$--$21\%$ on the training pools). We therefore take a methodological position instead of
quoting a number: \textbf{there is no reliable way to score a method inside the occlusion region
of this panel}. That matches the identifiability ceiling of Section~\ref{sec:diag} and is why our
supervision-placement study is run during \emph{training} rather than from per-region scores.

\subsection{Does fine-tuning on our data help?}
\label{sec:ft}

The practical question a dataset must answer is whether it transfers. We fine-tune the same
backbone (RDNet with Mona adapters, frozen FocalNet, $65.95$~M trainable) on our real pairs
and evaluate on the frozen 15-scene panel; the recipe is 6000 iterations, batch 6,
$\mathrm{lr}=10^{-4}$, mixed with online synthesis and open-source pairs at $30/40/30\%$.

\begin{table*}[t]
\centering
\caption{Fine-tuning on our data, evaluated on the frozen 15-scene panel. Which months the
pool contains matters far more than whether those months contain occlusion regions.
$^\ddagger$ rows apply a post-hoc per-image affine (the photometric upper bound of
Table~\ref{tab:photometric}).}
\label{tab:ft}
\small
\setlength{\tabcolsep}{6pt}
\begin{tabular}{@{}lccccc@{}}
\toprule
Training data & occlusion in pool & \#real pairs & PSNR$\uparrow$ & hole & non-hole \\
\midrule
identity (no training) & --- & 0 & 18.44 & 15.17 & 18.97 \\
official RDNet, zero-shot & --- & $\sim$10k (SIR2) & 20.05 & 19.25 & 20.19 \\
RRW, zero-shot & --- & 14{,}952 & 20.50 & 19.21 & 20.68 \\
\midrule
\multicolumn{6}{@{}l}{\emph{earlier recipe, single seed, small pools (934$\times$955 crop)}}\\
ours: may & \ding{55} & 14 & 20.69 & 18.91 & 20.98 \\
ours: apr & \ding{51} & 16 & 19.26 & 16.26 & 19.72 \\
ours: apr $+$ may & mixed & 30 & 20.63 & 18.74 & 20.92 \\
ours: may $+$ jun & mixed & 27 & 20.80 & \textbf{21.44} & 20.76 \\
ours: apr $+$ jun & \ding{51}\ding{51} & 29 & 21.32 & 20.21 & 21.47 \\
\midrule
\multicolumn{6}{@{}l}{\emph{current recipe, three seeds, full pools}}\\
ours: apr $+$ jun & \ding{51}\ding{51} & 170 & 20.97 & 18.20 & 21.39 \\
ours: apr $+$ jun$^\ddagger$ & \ding{51}\ding{51} & 170 & 22.13 & --- & --- \\
\textbf{ours: may $+$ jun} & mixed & \textbf{174} & \textbf{22.70} & \textbf{20.67}
  & \textbf{23.03} \\
ours: may $+$ jun$^\ddagger$ & mixed & 174 & \textbf{23.14} & --- & --- \\
\bottomrule
\end{tabular}

\vspace{2pt}
\footnotesize \emph{hole / non-hole} is the v13 partition ($r=R/T$ gate $+$ Otsu) computed by the
evaluation script: a \emph{reflection} mask, not the training-time occlusion region $\Omega$. The
sign of the gap depends on this partition, which is why we never use it as evidence on its own.
All rows share the same recipe (6000 iterations, batch 6, $\lambda_\Omega=3$); the two
three-seed rows differ \emph{only} in which months the pool is drawn from
(apr $+$ jun: 126$+$44; may $+$ jun: 130$+$44), and their seeds span $0.33$--$0.46$~dB.
\end{table*}

\textbf{The region term itself is worth $3$--$4$~dB.} Before asking \emph{where} supervision
should be placed, we measure what the region term is worth at all, by removing it entirely
($\lambda_{om}=0$) from the same recipe. On the april$+$june pool it costs $3.19$~dB
($20.96\!\to\!17.77$) and on may$+$june $4.38$~dB ($22.71\!\to\!18.33$); in both cases the
fine-tuned model ends up \emph{below} the identity baseline ($18.44$), so without this term
fine-tuning on our data is a net loss. This is the largest effect we measure. It also explains
the pool gap: it shrinks from $1.73$~dB to $0.56$~dB once the term is removed, i.e.\ most of the
may$+$june advantage only exists while the term is present.

\textbf{Reading.} Pool choice dominates everything we varied in this section. The
May-dominant pool reaches $22.70$~dB, i.e.\ $\mathbf{+2.22}$ over the strongest public
method (per-scene 95\% CI $[+1.73,+2.72]$, winning on \textbf{15/15} scenes) and $+2.68$ over
the official model, while the April-dominant pool of the same size and recipe reaches only
$20.97$ ($+0.47$, CI $[-0.03,+0.93]$, 11/15 --- not significant). Unlike every earlier arm,
the winning pool \emph{also} improves the hole region ($20.67$ vs.\ $19.21$ and $19.25$).

\textbf{How much of that is photometry?} Table~\ref{tab:photometric} supplies the photometric
upper bound: a post-hoc per-image affine lifts the April pool by $+1.21$~dB and the May pool
by only $+0.43$~dB, so roughly $0.8$~dB of the $1.77$~dB gap is the April pool's residual
brightness bias. \textbf{About $1.0$~dB survives that correction} ($23.14$ vs.\
$22.13$), so pool choice is \emph{not} purely photometric; we do not yet attribute the
remainder, and we report the raw numbers rather than the corrected ones throughout.

 \textbf{An analytically scale-invariant reconstruction does not help.} We also tried
replacing the reconstruction term by a per-sample affine-invariant one (closed-form $a,b$,
$a\!\ge\!0.2$, normalized by $\sqrt{\mathrm{Var}(T)}$). Two seeds give $20.46$ and $17.77$ against
the $20.96$ linear baseline --- no gain and a $2.7$~dB seed spread --- so the variant was dropped
and we report it as a negative result rather than a tuning detail.

\textbf{What this corrects in our earlier reading.} We previously reported that the ordering
follows \emph{whether the training months contain occlusion regions}. With three seeds and
matched recipes that is false: the pool dominated by the occlusion-free month is the better
one, by $1.7$~dB, and it is also better inside the hole. The variable that actually moves the
score is which acquisition month the pool represents, i.e.\ how close its photometry and scene
statistics are to the evaluation panel --- not the presence of occlusion supervision. We had
expected the opposite (hole supervision should be what a tinted-glass pool adds); the data say
otherwise, and we report the correction rather than the original story.

\textbf{The earlier 21.32 does not survive the controlled re-run} ($20.97$), confirming that
the ordering among the single-seed rows above is dominated by pool, crop and seed rather than
by whether a month contains occlusion regions.

\subsection{Where should supervision be placed?}
\label{sec:diag}

We then ask \emph{which part} of a region loss carries the signal. Identical arms differ only
in what the loss acts on (same pool of 82 human-confirmed pairs, three seeds per arm).

\begin{table*}[t]
\centering
\caption{Supervision placement. Means over three seeds; ``random'' is an area-matched disk
redrawn every iteration. Rows below the rule are single-seed screening runs.}
\label{tab:placement}
\small
\setlength{\tabcolsep}{6pt}
\begin{tabular}{@{}lccccc@{}}
\toprule
What the loss acts on & PSNR & seed sd & hole & non-hole & LapVar \\
\midrule
no region loss ($\lambda_\Omega=0$) & 17.76 & 0.16 & 14.44 & 18.29 & 34.5 \\
true $\Omega$ mask & 19.41 & \textbf{0.54} & 16.81 & 19.80 & 40.2 \\
random disk, redrawn & \textbf{19.90} & \textbf{0.16} & \textbf{17.42} & \textbf{20.26} & 36.1 \\
\midrule
whole image, no mask$^\dagger$ & 20.05 & --- & 17.42 & 20.43 & 36.0 \\
random disk, fixed per pair & 20.25 & --- & 17.70 & 20.62 & 37.2 \\
true mask $+$ true darkening$^\dagger$ & 20.06 & --- & 17.22 & 20.49 & --- \\
\bottomrule
\end{tabular}
\vspace{2pt}
\footnotesize $^\dagger$ single seed; screening evidence only.
\end{table*}

Three findings follow.

\textbf{The region need not be fixed in place either.} A disk fixed per pair scores $20.25\pm0.29$
over three seeds against $19.90\pm0.16$ for a disk redrawn every iteration --- a gap of $0.35$~dB,
below the $0.5$~dB threshold we pre-registered, so we report no significant difference. The
single-seed $20.48$ of an earlier run did not replicate, which is why all placement claims in
this paper rest on three seeds.

\textbf{(1) The loss works, but not because of the hole.} Removing it costs more than doing
nothing ($17.76$ versus $18.44$~dB for identity), so the term is load-bearing. Yet an
area-matched random disk reaches $+2.15$~dB over the no-loss arm while the \emph{true} mask
reaches $+1.66$~dB; the difference is $-0.49$~dB (1 of 3 seeds positive). The benefit is
``some region is under supervision'', not ``this region is the occlusion''.

\textbf{(2) Seed variance is itself informative.} The true-mask arm has a seed standard
deviation of $0.54$~dB while the random-disk and no-loss arms have $0.16$~dB. A mask that
coincides with real structure admits two nearly equivalent solutions (preserve it, or smooth
it) and which one is reached depends on initialization; a random disk does not. This is why
single-seed comparisons below $1$~dB are meaningless here.

\textbf{(3) What carries the signal is the normalization.} Replacing the masked relative error
by a \emph{whole-image} relative error already gives $20.05$~dB --- more than the true mask
--- identifying the denominator $\sum|\hat R|$ as the active ingredient. Consistently,
supplying the true in-region darkening removes the harm of the true mask ($20.06$~dB): the
mask hurt because it forced the model to guess an illumination level that one frame does not
contain.

\textbf{The boundary term cannot be judged.} Penalizing $|\nabla\hat T|>\alpha|\nabla\hat R|$
on $\partial\Omega$ gave $+0.72$~dB on one seed and $-0.37/-0.17$ on two more (mean
$+0.06$~dB, 1/3 positive), and occupies $0.06$--$0.5\%$ of the total loss. The support-area
bound of Eq.~\eqref{eq:bound} explains why: a 2-pixel band cannot move a global score by more
than $0.036$~dB, so the experiment was unidentifiable by construction.

\subsection{The reflection mask: better labels, no downstream gain}
\label{sec:gate}

Prior work obtains a spatial reflection prior from the gradient comparison $G_I>G_T$. On our
data that pseudo-label covers $48.1\%$ of the image and agrees with the true reflection energy
at only IoU $0.063$. Supervising the same head with the \emph{true} reflection layer instead
raises the mask to IoU $0.376$ (precision $0.739$, recall $0.489$), a $6\times$ improvement at
equal cost and with the same test-time interface (image only).

\textbf{But better labels do not yet buy accuracy.} Used as a per-pixel loss weight on the
same base arm, the true-reflection label gives $19.75$~dB while the original
gradient-comparison label gives $20.36$~dB, and both are below simply using no mask at all
($20.05$~dB, Table~\ref{tab:placement}). We therefore report the mask as a
\emph{diagnostic} --- it shows that the widely used pseudo-label is nearly uninformative ---
and not as a component with a demonstrated gain; closing that gap (e.g.\ as a conditioning
input rather than a weight) is future work.

\subsection{Photometric axis and an interventional test}
\label{sec:photom}

\begin{table*}[t]
\centering
\caption{Brightness is a large, non-identifiable component of the error. A closed-form
2-parameter affine per image, fitted post hoc, recovers most of the gap between the official
model and the trivial baseline. Net gain of the official model over that baseline: $+0.30$~dB
(apr), $-0.32$~dB (may).}
\label{tab:photometric}
\small
\setlength{\tabcolsep}{6pt}
\begin{tabular}{@{}lcccccc@{}}
\toprule
Test set & $n$ & identity, raw & identity, $+$affine & identity, $+$scalar
  & official, raw & official, $+$affine \\
\midrule
apr (\texttt{bak\_1026}, 16 frames) & 16 & 17.81 & \textbf{20.10} & 19.52
  & 17.42 & 20.40 \\
may & 14 & 23.32 & \textbf{25.77} & 25.61 & 23.71 & 25.44 \\
\bottomrule
\end{tabular}
\end{table*}

On the apr domain a two-parameter global affine applied to the \emph{identity} output reaches
$20.10$~dB while the official model gives $17.42$~dB; the net gain of the official model over
that trivial baseline is only $+0.30$~dB (and $-0.32$~dB on may). The bottleneck is therefore
photometric, not training effort --- what the non-identifiability argument predicts.

Over 35 fine-tuned arms (12-scene held-out panel, which excludes the three scenes used to build one training split), the systematic brightness bias is $6.66$ gray levels on average, and
a closed-form gain $g=\mu(T)/\mu(\hat T)$ recovers $+0.63$~dB on average ($\max$ $+1.89$~dB)
without changing the ranking. Since this benchmark's detection floor for brightness-only
effects is $1.19$--$1.31$~dB, the photometric error alone accounts for a measurable effect
size, and it is exactly the one global scale that a single frame cannot supply.

To test whether a model really uses a photometric condition, we perturb the conditioning
scalar $c$ and measure
$\mathrm{ControlGain}=\bigl(\partial\mu(\hat T)/\partial c\bigr)/\mu(I)$, requiring
$\ge0.50$. A frozen state-of-the-art model scores $0.02$, an absorption-aware network $0.086$,
and a model trained with a counterfactually sampled scalar $0.244$; the uninformative-condition
control scores $0.020$, a factor $12.4$ below the counterfactual arm.
